\documentclass[aspectratio=169,11pt]{beamer}
% Shared look for every deck in this repo: Metropolis with a progress bar under each
% frame title. Mirrors sourdough-science/slides/preamble.tex. Build with Tectonic:
% `make` in this directory.

\usetheme[progressbar=frametitle,numbering=fraction,sectionpage=progressbar,block=fill]{metropolis}

% Metropolis looks Fira Sans up by system name; load it from the TeX bundle instead,
% Light for body text as Metropolis intends, with lining figures for data.
\usepackage[sfdefault,light,lf]{FiraSans}
\usepackage{booktabs,tabularx,array}
\usepackage{tikz}
\usetikzlibrary{arrows.meta,positioning,shapes.geometric,calc}
\usepackage{pgfplots}
\pgfplotsset{compat=newest}

% Metropolis's palette, with this repo's own colour names mapped onto it.
\definecolor{mDarkTeal}{HTML}{23373B}
\definecolor{mLightBrown}{HTML}{EB811B}
\colorlet{ink}{mDarkTeal}
\definecolor{panel}{HTML}{EDF1F2}
\colorlet{accent}{mLightBrown}
\definecolor{flag}{HTML}{B03A2E}
\definecolor{ok}{HTML}{2E7D4F}
\definecolor{muted}{HTML}{7A7F80}

% Top-align frames, so the \vfill in \source puts the source line at the bottom.
% A slimmer title bar and narrower margins than Metropolis's defaults leave room
% for the denser slides.
\makeatletter
\beamer@centeredfalse
\setlength{\metropolis@frametitle@padding}{1.4ex}
\setlength{\metropolis@progressinheadfoot@linewidth}{2pt}
\setlength{\metropolis@progressonsectionpage@linewidth}{1.5pt}
\setlength{\metropolis@titleseparator@linewidth}{1pt}
% Metropolis's title page is \paperheight tall, which overflows a top-aligned frame.
\setbeamertemplate{title page}{%
  \begin{minipage}[b][\textheight]{\textwidth}
    \vfill
    \usebeamertemplate*{title}
    \ifx\insertsubtitle\@empty\else\usebeamertemplate*{subtitle}\fi
    \usebeamertemplate*{title separator}
    \usebeamertemplate*{author}
    \usebeamertemplate*{institute}
    \usebeamertemplate*{date}
    \vfill
  \end{minipage}}
\makeatother
\setbeamersize{text margin left=1.2em,text margin right=1.2em}

\setbeamertemplate{frame footer}{\insertshorttitle}
\setbeamertemplate{footline}{%
  \begin{beamercolorbox}[wd=\paperwidth,ht=2.6ex,dp=1.2ex,leftskip=1.2em,rightskip=1.2em]{footline}%
    \usebeamerfont{page number in head/foot}%
    \usebeamertemplate*{frame footer}\hfill\usebeamertemplate*{frame numbering}%
  \end{beamercolorbox}}
\setbeamerfont{page number in head/foot}{size=\tiny}
\setbeamercolor{footline}{fg=muted}

% The conclusion of a slide, in one sentence.
\newcommand{\reading}[1]{%
  \par\medskip
  \colorbox{panel}{\parbox{\dimexpr\linewidth-2\fboxsep}{\small\textbf{\color{ink}What it means:} #1}}}

% Small source line, placed at the bottom of a frame.
\newcommand{\source}[1]{%
  \vfill{\color{muted}\tiny Source: #1}}

% Status tags. Filled, so they read on the dark title bar as well as on the body.
\newcommand{\flagtag}[2]{\mbox{\setlength{\fboxsep}{1.2pt}\colorbox{#1}{\color{white}\scriptsize\bfseries\,#2\,}}}
\newcommand{\donetag}{\flagtag{ok}{done}}
\newcommand{\presenttag}{\flagtag{ok}{present}}
\newcommand{\pendingtag}{\flagtag{accent}{pending}}
\newcommand{\ruledout}{\flagtag{muted}{ruled out}}
\newcommand{\phantomtag}{\flagtag{flag}{Code 45}}
\newcommand{\hyp}{\flagtag{flag}{hypothesis}}
\setbeamertemplate{frametitle continuation}[from second][(cont.)]

% Nodes for TikZ diagrams: a normal box, and a box for a device that is missing.
\tikzset{
  box/.style={draw=ink,fill=panel,rounded corners=2pt,align=center,
    font=\scriptsize,inner sep=4pt,text width=#1},
  box/.default=2.6cm,
  gone/.style={draw=flag,dashed,fill=white,rounded corners=2pt,align=center,
    font=\scriptsize,inner sep=4pt,text width=#1},
  gone/.default=2.6cm,
  arrow/.style={-{Stealth[length=5pt]},draw=ink,thick},
  gonearrow/.style={-{Stealth[length=5pt]},draw=flag,dashed,thick},
}

\newcolumntype{L}{>{\raggedright\arraybackslash}X}
\newcommand{\muted}[1]{{\color{muted}#1}}

\author{ZRRBITE-PC hardware notes}
\institute{\texttt{C:/dev/hardware-notes}}


\title[RTX 3050 failure]{RTX 3050: gone from the PCIe bus}
\subtitle{What happened on ZRRBITE-PC, what the evidence says,\\
  and what to do next, including the BIOS update}
\date{Status as of 2026-11-01 evening: the card is back after the slot test}

\begin{document}

\begin{frame}[plain]\titlepage\end{frame}

% ---------------------------------------------------------------
\begin{frame}{Where things stand}
  \begin{itemize}\setlength{\itemsep}{0.2em}
    \item From 2026-10-30 16:41 to 2026-11-01 17:31 the RTX 3050 was not detected at power-on.
      Windows ran on the Intel UHD 770; the monitor on the onboard output.
    \item Card, HDMI audio function and the CPU's x16 root port were all \emph{phantom}
      (Device Manager Code 45): absent from the bus, not failing to start.
    \item 2026-11-01 morning: cold power cycle, card and 8-pin lead reseated. No change.
    \item 2026-11-01 17:31: card moved to a lower chipset slot, \textbf{detected} at Gen3 x4.
      17:37: back in the top slot, detected at Gen4 x8. Driver loaded; GPU 1 is back.
    \item 59 crash reports since 2026-05-04 remain unexplained. 621.01 did not cause the loss
      but lagged; 617.42 replaced it at 20:03, first games ran clean.
  \end{itemize}
  \reading{the card and the top slot both train again. The fault was in the link path and
    cleared with the slot change. Its cause is not identified, so it can come back.}
  \source{diagnosis/gpu-rtx3050-missing-2026-10.md}
\end{frame}

% ===============================================================
\section{The machine}

\begin{frame}{What is in the box}
  \footnotesize
  \renewcommand{\arraystretch}{1.08}
  \begin{tabularx}{\textwidth}{@{}>{\bfseries\raggedright}p{2.6cm}L@{}}
    \toprule
    Motherboard & Gigabyte Z690 AORUS ELITE DDR4 rev.\ 1.x, BIOS \textbf{F5} (2021-12-17, launch era) \\
    CPU & Intel Core i7-12700K; UHD Graphics 770 integrated, now driving the display \\
    Graphics card & Gigabyte GeForce RTX 3050 8\,GB, PCI 10DE:2582, subsystem 1458:4101, rev.\ A1.
      Top slot (CPU PCIe x16); the card is natively x8; 130\,W, one 8-pin lead \\
    Memory & 2\,$\times$\,16\,GB DDR4 (G.Skill F4-4000C18), running at 2666\,MT/s, XMP off \\
    Storage & Samsung 970 EVO Plus 500\,GB, Samsung 980 PRO 1\,TB (Windows) \\
    Network & Realtek RTL8125 2.5GbE on a chipset root port (bus 5, Gen2 x1) \\
    Power supply & \textbf{unknown}: label still to be read \\
    OS & Windows 11 Pro for Workstations, build 26200, installed 2026-02-14.
      Hibernation and Fast Startup off since 2026-09-05, so Shut down is a true power-off \\
    Monitor & Lenovo P34w-20, 3440\,$\times$\,1440 \\
    \bottomrule
  \end{tabularx}
  \source{hardware/inventory.md, scripts/Get-HardwareInventory.ps1}
\end{frame}

% ===============================================================
\section{What happened}

\begin{frame}{Timeline}
  \scriptsize
  \renewcommand{\arraystretch}{1.0}
  \begin{tabularx}{\textwidth}{@{}>{\raggedright}p{3.1cm}L@{}}
    \toprule
    \textbf{When} & \textbf{Event} \\
    \midrule
    2026-05-04 & First GPU timeout dump (0x141). Fifteen more in May. \\
    2026-06 to 2026-10 & Escalation: black-screen dumps, graphics watchdog dumps, seven blue screens \\
    2026-10-23 04:43 & Driver 610.88 arrives via Windows Update; blue screen 0x7E in nvlddmkm the same day \\
    2026-10-29 13:53 & Blue screen 0x50 in nvlddmkm \\
    2026-10-29 17:48 & Driver 621.01 installed via NVIDIA App. Clean install, no firmware update. Card runs 22\,h \\
    2026-10-30 16:12:22 & System hangs. Last log write; no crash dump \\
    2026-10-30 16:41:18 & Power button held to force off (Kernel-Power 41 carries the timestamp). Boot at 16:41:44: card, audio, root port absent \\
    2026-10-30 16:43:13 & Second power-button hold; boot at 16:44:00, same result. Monitor moved to onboard HDMI \\
    2026-10-31 & Software diagnosis, SSH set up as a lifeline, vendor report drafted \\
    2026-11-01 08:48 to 09:00 & Cold power cycle, card and 8-pin reseated, inspection clean. Two clean boots, card still absent \\
    \bottomrule
  \end{tabularx}
  \source{System log (Kernel-Power, Kernel-General), WER ReportArchive, NVIDIA setup log}
\end{frame}

% ---------------------------------------------------------------
\begin{frame}{The signature: three devices missing at once}
  \centering
  \scalebox{0.88}{\begin{tikzpicture}[node distance=0.25cm and 0.55cm]
    \node[box=2.5cm] (cpu) {\textbf{CPU}\\i7-12700K};
    \node[gone=3.1cm,right=of cpu] (peg)
      {\textbf{PEG10 root port}\\8086:460D, bus 0 dev 1\\\phantomtag};
    \node[gone=2.7cm,right=of peg] (slot)
      {\textbf{PCIEX16 slot}\\CPU PCIe 5.0 x16\\\muted{no link trained}};
    \node[gone=3.2cm,anchor=west] (gpu) at ($(slot.east)+(0.55,0.78)$)
      {\textbf{RTX 3050} (fn 0)\\10DE:2582\\\phantomtag};
    \node[gone=3.2cm,anchor=west] (aud) at ($(slot.east)+(0.55,-0.78)$)
      {\textbf{HDMI audio} (fn 1)\\10DE:2291\\\phantomtag};
    \node[box=2.5cm,below=0.25cm of cpu] (igpu) {\textbf{UHD 770}\\8086:4680\\\presenttag};
    \node[box=2.5cm,below=0.25cm of igpu] (pch) {\textbf{Z690 chipset}};
    \node[box=3.1cm,right=of pch] (rp) {\textbf{Root port}\\8086:7ABA, bus 0 dev 28 fn 2\\\presenttag};
    \node[box=2.7cm,right=of rp] (nic) {\textbf{RTL8125} 2.5GbE\\\muted{corrected errors\\at every boot}};
    \draw[gonearrow] (cpu) -- (peg);
    \draw[gonearrow] (peg) -- (slot);
    \draw[gonearrow] (slot.east) -- (gpu.west);
    \draw[gonearrow] (slot.east) -- (aud.west);
    \draw[arrow] (cpu) -- (igpu);
    \draw[arrow] (igpu) -- (pch);
    \draw[arrow] (pch) -- (rp);
    \draw[arrow] (rp) -- (nic);
  \end{tikzpicture}}
  \par\smallskip\raggedright\scriptsize
  The root port belongs to the CPU, not the card; the firmware exposes it only when something trains
  a link in that slot at POST. Three phantoms, no removal timestamps: \emph{card not detected at POST}.
  \source{scripts/Get-GpuState.ps1 on 2026-11-01; last recorded link on 2026-10-29 was Gen4 x8, no AER errors}
\end{frame}

% ---------------------------------------------------------------
\begin{frame}{Why it is not the driver}
  \begin{itemize}\setlength{\itemsep}{0.25em}
    \item The root port is enumerated by the BIOS before Windows loads. No driver runs yet.
    \item The three devices carry no removal timestamp: not disabled or hot-removed in a
      session, simply absent at the boots of 16:41 and 16:44 and at every boot since.
    \item 621.01 installed cleanly: no firmware or GOP update component, no reboot required,
      and the card ran normally for 22 hours afterwards.
    \item Faults span four driver packages over six months (595.79, 610.74, 610.88, 621.01).
    \item The 621.01 driver package and its display container service are still installed;
      \texttt{nvidia-smi} reports no device.
  \end{itemize}
  \reading{rolling back or reinstalling a driver cannot bring back a device the firmware does
    not see. Driver work only matters after the card reappears, and then as a DDU clean install.}
  \source{drivers/nvidia.md; NVIDIA LOG.setup.exe.log; Kernel-PnP and UserPnp events}
\end{frame}

% ---------------------------------------------------------------
\begin{frame}{Six months of warnings}
  \begin{columns}[T,onlytextwidth]
    \begin{column}{0.56\textwidth}
      \centering
      \begin{tikzpicture}
        \begin{axis}[ybar stacked,bar width=17pt,width=\linewidth,height=4.5cm,
            symbolic x coords={May,Jun,Jul,Aug,Sep,Oct},xtick=data,
            ymin=0,ymax=18,ylabel={crash reports per month},
            ylabel style={font=\scriptsize},tick label style={font=\scriptsize},
            axis lines*=left,enlarge x limits=0.12,ymajorgrids,grid style={muted!30},
            legend style={font=\tiny,at={(0.98,0.97)},anchor=north east,legend columns=1,draw=none,fill=none}]
          \addplot[ybar,fill=ink!75,draw=none] coordinates {(May,15) (Jun,12) (Jul,0) (Aug,0) (Sep,0) (Oct,0)};
          \addlegendentry{GPU timeout, 0x141/0x117}
          \addplot[ybar,fill=accent!85,draw=none] coordinates {(May,0) (Jun,2) (Jul,8) (Aug,0) (Sep,0) (Oct,2)};
          \addlegendentry{black screen, 0x1A8/0x1B8}
          \addplot[ybar,fill=muted!70,draw=none] coordinates {(May,0) (Jun,0) (Jul,3) (Aug,1) (Sep,6) (Oct,3)};
          \addlegendentry{watchdog, 0x193/0x1CC}
          \addplot[ybar,fill=flag,draw=none] coordinates {(May,0) (Jun,2) (Jul,0) (Aug,1) (Sep,2) (Oct,2)};
          \addlegendentry{blue screen}
        \end{axis}
      \end{tikzpicture}
    \end{column}
    \begin{column}{0.42\textwidth}
      \scriptsize
      59 kernel reports from 2026-05-04 to 2026-10-31, all in the graphics stack:
      52 live kernel events and 7 blue screens.
      \smallskip
      \begin{itemize}\setlength{\itemsep}{0.2em}
        \item \textbf{May to June:} recoverable GPU timeouts; the driver reset the GPU and carried on
        \item \textbf{July:} black-screen diagnostics; the display dropped out
        \item \textbf{August to October:} kernel crashes in GPU memory-management code
        \item \textbf{30 October:} the card stops enumerating
      \end{itemize}
      \smallskip
      The two October black-screen dumps (2026-10-31) are from the integrated GPU, after the card
      was gone.
    \end{column}
  \end{columns}
  \source{C:/ProgramData/Microsoft/Windows/WER/ReportArchive; scripts/Get-CrashHistory.ps1}
\end{frame}

% ---------------------------------------------------------------
\begin{frame}{The seven blue screens}
  \scriptsize
  \renewcommand{\arraystretch}{1.05}
  \begin{tabularx}{\textwidth}{@{}p{2.9cm}p{1.5cm}L@{}}
    \toprule
    \textbf{When} & \textbf{Code} & \textbf{Where it crashed} \\
    \midrule
    2026-10-29 13:53 & 0x50 & nvlddmkm!\_clientUnmapInterBackRefMappings (NVIDIA kernel driver) \\
    2026-10-23 18:35 & 0x7E & nvlddmkm!refRemoveInterMapping \\
    2026-09-26 09:28 & 0x7E & dxgmms2!VIDMM\_GLOBAL::UnlockAllocation (DirectX video memory manager) \\
    2026-09-24 00:06 & 0x7E & dxgmms2!VIDMM\_GLOBAL::MakeOneVirtualAddressRangeNotResident \\
    2026-08-24 12:20 & 0xA  & nvlddmkm!CNvLAllocReference::unlinkFromAllLists \\
    2026-06-06 23:07 & 0xA  & nt!MiCrashdumpRemovePte, while writing a crash dump \\
    2026-06-06 12:58 & 0x133 & DPC watchdog, nvlddmkm!CNvLAdapter::lockallSubmissions \\
    \bottomrule
  \end{tabularx}
  \reading{five of seven are inside the NVIDIA driver or DirectX's handling of the card's memory.
    The signatures change from release to release; the component that fails does not.}
  \source{WER ReportArchive folder names and Report.wer buckets; Application log, Windows Error Reporting 1001}
\end{frame}

% ===============================================================
\section{What the hardware is telling us}

\begin{frame}{Power and fans}
  \begin{columns}[T,onlytextwidth]
    \begin{column}{0.44\textwidth}
      \centering
      \begin{tikzpicture}
        \node[box=1.4cm] (psu) {\textbf{PSU}\\model unknown};
        \node[box=1.6cm,right=0.45cm of psu,yshift=0.7cm] (slot) {\textbf{Slot}\\up to 75\,W};
        \node[box=1.6cm,right=0.45cm of psu,yshift=-0.7cm] (aux) {\textbf{8-pin lead}\\up to 150\,W};
        \node[box=1.6cm,right=0.45cm of slot,yshift=-0.7cm] (card) {\textbf{RTX 3050}\\130\,W rated};
        \draw[arrow] (psu.east) -- (slot.west);
        \draw[arrow] (psu.east) -- (aux.west);
        \draw[arrow] (slot.east) -- (card.west);
        \draw[arrow] (aux.east) -- (card.west);
      \end{tikzpicture}
      \par\medskip\scriptsize\raggedright
      No software view of the card's power exists while it is off the bus. The 12\,V rails can be
      checked at a spare connector with a multimeter (ATX tolerance $\pm$5\,\%) or a PSU tester;
      HWMonitor only sees board sensors.
    \end{column}
    \begin{column}{0.54\textwidth}
      \footnotesize
      \textbf{Observed while the card is undetected}
      \begin{itemize}\setlength{\itemsep}{0.3em}
        \item Both fans spin for a few seconds at power-on, then stop: the firmware's fan
          self-test, then fan-stop because the chip is cold. A healthy card does the same.
        \item One fan occasionally restarts at idle, and the heatsink gets hot with no load:
          without a driver the GPU sits at boot clocks and burns tens of watts.
        \item A card with no 12\,V, or a dead core, would show neither.
      \end{itemize}
    \end{column}
  \end{columns}
  \reading{the card is powered and alive at the firmware level. The fault is in PCIe link
    training, on the card or on the board's x16 port.}
  \source{owner observations 2026-10-31 and 2026-11-01; hardware/motherboard.md, PSU measurement notes}
\end{frame}

% ---------------------------------------------------------------
\begin{frame}{What has been ruled out}
  \scriptsize
  \renewcommand{\arraystretch}{1.0}
  \begin{tabularx}{\textwidth}{@{}p{3.3cm}Lp{1.7cm}@{}}
    \toprule
    \textbf{Hypothesis} & \textbf{Test} & \textbf{Verdict} \\
    \midrule
    Driver 621.01 broke it & Root port is hidden by firmware before Windows loads; clean install log; 22\,h of normal operation & \ruledout \\
    PSU dropped out & Kernel-Power 41 events carry power-button timestamps; no power-loss or crash shutdowns in the log & \ruledout \\
    Card or cable loose & Reseated card and 8-pin on 2026-11-01; connectors and board inspected; no change & \ruledout \\
    8-pin not delivering & Fans self-test at power-on; heatsink warms. A card without auxiliary power would refuse to start & \ruledout \\
    Link degraded before failure & Last recorded link Gen4 x8; no AER errors on the card; WHEA errors only on the network card's port & \ruledout \\
    Monitor or Windows install & Integrated GPU drives the same monitor at full resolution & \ruledout \\
    Card's PCIe interface (PHY or BGA joints) & Card in a chipset slot, then another PC & \hyp \\
    Board's CPU x16 port, slot or CPU lanes & Card in a chipset slot; BIOS link-speed test; F35 flash & \hyp \\
    \bottomrule
  \end{tabularx}
  \source{diagnosis/gpu-rtx3050-missing-2026-10.md, Evidence and Ruled out sections}
\end{frame}

% ---------------------------------------------------------------
\begin{frame}{What remains: card or board}
  \begin{columns}[T,onlytextwidth]
    \begin{column}{0.4\textwidth}
      \footnotesize
      \textbf{First, without opening the case}
      \begin{enumerate}\setlength{\itemsep}{0.3em}
        \item Enter setup (Del). Look for the card anywhere in the device listing.
        \item Top slot link speed: Gen3, then Gen1.
        \item Load Optimized Defaults (F7), Initial Display Output = IGFX, save (F10).
        \item Boot, run \texttt{scripts/Get-GpuState.ps1}.
      \end{enumerate}
      \medskip
      \textbf{Run 2026-11-01:} the card appeared in both slots. Neither branch applied;
      the link was marginal, not dead.
    \end{column}
    \begin{column}{0.58\textwidth}
      \centering
      \begin{tikzpicture}[node distance=0.45cm and 0.3cm]
        \node[box=4.4cm] (q) {\textbf{Move the card to a lower slot}\\chipset, x4 electrical; boot};
        \node[box=3.0cm,below left=0.5cm and -1.5cm of q] (yes)
          {\textbf{Card appears}\\card is fine; CPU x16 port, slot or CPU lanes at fault};
        \node[box=3.0cm,below right=0.5cm and -1.5cm of q] (no)
          {\textbf{Still absent}\\card does not train PCIe at all};
        \node[box=3.0cm,below=of yes] (flash)
          {Flash F35, retest the top slot.\\Still absent: board or CPU under warranty};
        \node[box=3.0cm,below=of no] (rma)
          {Confirm in another PC if possible.\\Vendor report for the card};
        \draw[arrow] (q.south) -- (yes.north);
        \draw[arrow] (q.south) -- (no.north);
        \draw[arrow] (yes) -- (flash);
        \draw[arrow] (no) -- (rma);
      \end{tikzpicture}
    \end{column}
  \end{columns}
  \source{diagnosis note, Procedure steps 3 and 4; hardware/motherboard.md for the slot layout}
\end{frame}

% ---------------------------------------------------------------
\begin{frame}{The procedure and where we are}
  \scriptsize
  \renewcommand{\arraystretch}{1.05}
  \begin{tabularx}{\textwidth}{@{}p{0.5cm}Lp{2.8cm}@{}}
    \toprule
    & \textbf{Step} & \textbf{Status} \\
    \midrule
    1 & Cold power cycle: PSU off, power button held, wait, power on, watch the fans & \donetag\ 2026-11-01 \\
    2 & Reseat the card and the 8-pin lead; inspect contacts, slot, board and pins; read the PSU label & \donetag\ 2026-11-01, PSU label pending \\
    3 & BIOS settings: link speed Gen3 then Gen1, Optimized Defaults. Overtaken by step 4 & \flagtag{muted}{skipped} \\
    4 & Card or board: card in a chipset slot. Result: detected there, and again in the top slot & \donetag\ 2026-11-01 \\
    5 & If the card returns: block automatic driver installs, watch for a week, DDU clean install & \flagtag{accent}{active} \\
    6 & If it does not, or vanishes again: finish the vendor report (PSU, serial, purchase) and submit & \flagtag{muted}{on hold} \\
    \bottomrule
  \end{tabularx}
  \reading{steps 1 and 2 told us the card was powered. Step 4 brought it back without saying why.
    Step 5 is where we are.}
  \source{diagnosis/gpu-rtx3050-missing-2026-10.md, Procedure and Log}
\end{frame}

% ===============================================================
\section{Updating the BIOS}

\begin{frame}{Why F5 to F35, and where it fits}
  \begin{columns}[T,onlytextwidth]
    \begin{column}{0.5\textwidth}
      \footnotesize
      \textbf{What changed between F5 (2021-12) and F35}
      \begin{itemize}\setlength{\itemsep}{0.3em}
        \item \textbf{F23:} fix for a blank POST display with GeForce RTX 30-series cards
        \item \textbf{F32:} Intel CSME firmware update. One-way: a BIOS older than F32 cannot be
          flashed back afterwards
        \item Microcode, memory compatibility and stability updates throughout
      \end{itemize}
      \medskip
      Gigabyte allows the direct jump. Re-read the release notes on the support page before
      flashing and confirm nothing says ``update to Fxx first''.
    \end{column}
    \begin{column}{0.48\textwidth}
      \footnotesize
      \textbf{The image, already downloaded and checked}
      \begin{itemize}\setlength{\itemsep}{0.3em}
        \item \texttt{C:/dev/mobo/bios/\ldots/}\\\texttt{Z690AORUSELITEDDR4.F35}, 32\,MiB
        \item SHA-256 \texttt{55d90960\ldots{}3bc3ced00}
        \item Contains the board string \texttt{Z690 AORUS ELITE DDR4}
        \item Not committed to the repo
      \end{itemize}
      \medskip
      \textbf{Where it fits}\\
      The card ran for years on F5, so the flash alone is unlikely to revive it. Right once the
      slot test points at the board; harmless either way.
    \end{column}
  \end{columns}
  \source{bios/update-bios.md; Gigabyte support page for Z690 AORUS ELITE DDR4 rev.\ 1.x}
\end{frame}

% ---------------------------------------------------------------
\begin{frame}{Before you flash}
  \footnotesize
  \begin{enumerate}\setlength{\itemsep}{0.15em}
    \item \textbf{BitLocker.} In an elevated terminal (Win+X, Terminal (Admin)):
      \texttt{manage-bde -status C:}\\
      If protection is on, have the recovery key to hand (Microsoft account, Devices) and suspend
      it for one reboot: \texttt{manage-bde -protectors -disable C: -RebootCount 1}\\
      A BIOS change alters the TPM measurements; without this Windows may demand the key.
    \item \textbf{Record the non-default settings.} As of 2026-10-31: boot order (980 PRO) and
      Initial Display Output. XMP is off.
    \item \textbf{Power.} Steady mains, no storm. A power loss mid-write leaves the board
      unbootable; Q-Flash Plus is the recovery path.
    \item \textbf{Second device} with these notes; the machine is offline for a while.
    \item \textbf{USB stick}, FAT32, otherwise empty. Copy \texttt{Z690AORUSELITEDDR4.F35} to its root.
  \end{enumerate}
  \reading{do not flash from Windows (GIGABYTE Control Center, @BIOS), and never use a file for the
    DDR5 board or rev.\ 2.x. The file name must say DDR4.}
\end{frame}

% ---------------------------------------------------------------
\begin{frame}{Q-Flash, step by step}
  \begin{columns}[T,onlytextwidth]
    \begin{column}{0.5\textwidth}
      \scriptsize
      \begin{enumerate}\setlength{\itemsep}{0.2em}
        \item Shut down. Stick into a \textbf{rear} USB port, not a front-panel hub.
        \item Power on. Press \textbf{End} at the Gigabyte logo for Q-Flash, or \textbf{Del} for
          setup and then \textbf{F8}.
        \item \textbf{Update BIOS}, choose the stick, choose \texttt{Z690AORUSELITEDDR4.F35}, confirm.
        \item \textbf{Wait.} The screen may go dark and the machine may restart more than once
          (CSME update). Touch nothing until setup or Windows appears.
      \end{enumerate}
    \end{column}
    \begin{column}{0.48\textwidth}
      \scriptsize
      \begin{enumerate}\setlength{\itemsep}{0.2em}\setcounter{enumi}{4}
        \item Enter setup. \textbf{Load Optimized Defaults} (F7), \textbf{Save \& Exit} (F10).
        \item Enter setup again. Re-apply boot order and Initial Display Output. Leave XMP off
          unless you decide otherwise, and record it if you do.
        \item Boot Windows. Enter the BitLocker key once if asked. If you suspended protection,
          re-enable: \texttt{manage-bde -protectors -enable C:}
        \item Verify and record:\\
          \texttt{Get-CimInstance Win32\_BIOS | Select SMBIOSBIOSVersion, ReleaseDate}\\
          then update \texttt{hardware/inventory.md} and the log in \texttt{bios/update-bios.md}.
      \end{enumerate}
    \end{column}
  \end{columns}
  \reading{about fifteen minutes. The only step with a failure mode is the wait in step 4:
    a restart during the flash is normal, a power cut is not.}
  \source{bios/update-bios.md}
\end{frame}

% ---------------------------------------------------------------
\begin{frame}{If it will not POST afterwards: Q-Flash Plus}
  \begin{itemize}\setlength{\itemsep}{0.25em}
    \item This board family can flash from USB without a working BIOS.
    \item Rename the image to \texttt{GIGABYTE.bin}, put it alone on a FAT32 stick.
    \item Plug the stick into the rear port marked for Q-Flash Plus (check the manual for the exact port).
    \item System off, PSU switched on. Press the Q-Flash Plus button on the rear panel.
    \item Wait for the LED to stop blinking, several minutes. Then power on normally.
  \end{itemize}
  \reading{a failed flash is recoverable on this board: an inconvenience, not a dead motherboard.}
  \source{bios/update-bios.md; Gigabyte Z690 AORUS ELITE DDR4 manual}
\end{frame}

% ===============================================================
\section{Next}

\begin{frame}{Next steps}
  \begin{enumerate}\setlength{\itemsep}{0.15em}
    \item \textbf{Leave the card in the top slot.} Decline driver offers from NVIDIA App and Windows Update.
    \item \textbf{Watch it:} \texttt{scripts/Get-GpuState.ps1} after each boot for a week. Code 45
      again, or a link below Gen4 x8 under load, means the fault is back.
    \item \textbf{Crash history:} \texttt{scripts/Get-CrashHistory.ps1 -Since 2026-11-01} weekly.
      If the marginal link caused the May to October timeouts, they stop now.
    \item \textbf{Driver:} 617.42 since 2026-11-01 evening; 621.01 lagged. First games clean. DDU in reserve.
    \item \textbf{Flash F35} at leisure. Optional now; worth doing before any new card.
    \item \textbf{Record} the PSU label and the lower slot used; the vendor report stays a draft.
  \end{enumerate}
  \reading{every result goes in the Log table of the diagnosis note, dated. The vendor report
    is built from that table.}
\end{frame}

% ---------------------------------------------------------------
\begin{frame}{Sources and tools}
  \footnotesize
  \begin{itemize}\setlength{\itemsep}{0.3em}
    \item \texttt{diagnosis/gpu-rtx3050-missing-2026-10.md}: symptoms, hypothesis, evidence, ruled-out table, procedure, log
    \item \texttt{hardware/inventory.md}, \texttt{hardware/motherboard.md}: what is in the box; slots, firmware behaviour, rail tolerances
    \item \texttt{bios/update-bios.md}: the Q-Flash procedure and its log
    \item \texttt{drivers/nvidia.md}: driver versions, cached installers, clean-install plan
    \item \texttt{scripts/Get-GpuState.ps1}: presence, problem code, last arrival, PCIe link of the card and root port
    \item \texttt{scripts/Get-CrashHistory.ps1}: Kernel-Power 41 decoded, WHEA errors, WER kernel report history
    \item \texttt{scripts/Get-HardwareInventory.ps1}: board, BIOS, CPU, RAM, disks, video
    \item Vendor report (Claude Doc): \url{https://claude.ai/artifact/Vq7LVe2ABxcKtUVYbMFStL}
    \item Gigabyte support: \url{https://www.gigabyte.com/Motherboard/Z690-AORUS-ELITE-DDR4-rev-1x/support}
  \end{itemize}
  \source{all data collected on the affected machine, 2026-10-31 and 2026-11-01, without elevation}
\end{frame}

\end{document}
